% Source proofs adapted from pancyclic-paper.tex, Section6.
% Original source is unchanged. This file makes the new draft self-contained.
\section{Cycle-space, walk, and spectral estimates}
\label{sec:polynomial-foundations}

A region is a nonempty connected vertex set.
For a region $U$, let $\delta U$ be its edge boundary.
Write $\mathcal O_U$ for the event that a uniform even subgraph avoids $\delta U$.
For disjoint regions $U,V$, their links are the edges between them and the components of $G-U-V$ adjacent to both.
The number of links is $\ell(U,V)$.
The following proofs are adapted from Section6 of the existing manuscript.
They are included to expose every dependency of the polynomial argument.

\begin{lemma}[Link formula]\label{lem:link}
Let $A$ and $B$ be disjoint regions of a connected graph $G$. Then
$\PP(\cO_A)=2^{\comp(G-A)-|\delta A|}\ge2^{-|\delta A|}$ and
\[
  \PP(\cO_A\cap\cO_B)=2^{\ell(A,B)-1}\,\PP(\cO_A)\,\PP(\cO_B).
\]
Conditionally on $\cO_A$, the set $\xi\cap E(G[A])$ is uniform on $\Z(G[A])$; conditionally on
$\cO_A\cap\cO_B$, the sets $\xi\cap E(G[A])$ and $\xi\cap E(G[B])$ are independent and uniform on
$\Z(G[A])$ and $\Z(G[B])$.
\end{lemma}

\begin{proof}
The even subgraphs that avoid $\delta A$ are the sums of an even subgraph of $G[A]$ and an even
subgraph of $G-A$. So $\PP(\cO_A)=2^{\beta(G[A])+\beta(G-A)-\beta(G)}$, and counting edges and
vertices shows that the exponent is $\comp(G-A)-|\delta A|$. In the same way
\[
  \PP(\cO_A\cap\cO_B)=2^{\beta(G[A])+\beta(G[B])+\beta(G-A-B)-\beta(G)}
     =2^{\comp(G-A-B)+1-|\delta A|-|\delta B|+e(A,B)},
\]
where $e(A,B)$ is the number of edges between $A$ and $B$. Let $a$, $b$ and $k$ be the numbers of
components of $G-A-B$ that are joined to $A$ only, to $B$ only, and to both. Then
$\comp(G-A-B)=a+b+k$, while $\comp(G-A)=a+1$ and $\comp(G-B)=b+1$, and the formula follows since
$\ell(A,B)=k+e(A,B)$. The conditional laws are read off from the two direct sum decompositions.
\end{proof}

So two cut events are independent when there is one link. They are correlated by a factor of
exactly two when there are two links. The correlation is stronger only for pairs with at least
three links, and the next lemma says that such pairs are rare.

\begin{lemma}[Exceptional partners]\label{lem:exc}
Let $A$ be a region of a connected graph $G$, and let $B_1,\dots,B_k$ be pairwise disjoint regions,
disjoint from $A$, with $\ell(A,B_i)\ge3$ for all $i$. Then $k\le|\delta A|$.
\end{lemma}

\begin{proof}
Contract $A$ to a vertex $a$ and each $B_i$ to a vertex $b_i$, and delete the loops that arise.
This changes neither $d=|\delta A|$ nor the numbers $\ell(A,B_i)$, because the sets $B_j$ with
$j\ne i$ are connected and are contracted inside components of $G-(A\cup B_i)$. Let $F$ be a spanning forest of
$G-a$, and for every edge $e=au$ at $a$ attach a new leaf $\tau_e$ to $u$. Call the resulting forest
$F^{+}$ and its $d$ new leaves \emph{terminals}.

Fix $i$. Every link between $a$ and $b_i$ yields a path of $F^{+}$ from $b_i$ to a terminal. For an
edge $e=ab_i$ it is the edge $b_i\tau_e$. For a component $Y$ of $G-\{a,b_i\}$ joined to both, choose
an edge $e=au$ with $u\in Y$ and take the path of $F^{+}$ from $b_i$ to $\tau_e$; after its first
vertex it stays in $Y\cup\{\tau_e\}$. Paths that come from different links therefore start with
different edges at $b_i$, and the union of any two of them is a path between two terminals. Let
$F^{\circ}\subseteq F^{+}$ be the union of all paths between terminals. As $\ell(a,b_i)\ge3$, the
vertex $b_i$ has degree at least $3$ in $F^{\circ}$. All leaves of $F^{\circ}$ are terminals, and a
forest with at most $d$ leaves has at most $d-2$ vertices of degree at least $3$. Hence $k\le d-2$ if $k\ge1$.
\end{proof}


\subsection{Walks and normalized traces}
Let $J$ be a simple graph of maximum degree three and girth greater than $D\ge3$.
Put $s=\lfloor(D-1)/2\rfloor$ and let $L>D$.
Let $N_\ell(u,w)$ count non-backtracking walks of length $\ell$ from $u$ to $w$.
Write $v=|V(J)|$ and $d_J=\sum_u(3-\deg_Ju)$.
For the spectral supply estimate below, assume $d_J/v<1$.

\begin{lemma}\label{lem:walks}
If $\ell>s$, then $N_\ell(u,w)\le3\cdot2^{\ell-s-1}$ for all $u,w$. Consequently, for every vertex
$u$ of $J$, the cycles $C$ of $J$ through $u$ with $|C|\le L$ satisfy $\sum_C2^{-|C|}\le L\,2^{-s}$.
\end{lemma}

\begin{proof}
There are at most $3\cdot2^{\ell-s-1}$ ways to choose the first $\ell-s$ steps. The last $s$ steps
form a non-backtracking walk of length $s$ between two given vertices. Such a walk is a path,
because $s$ is less than the girth, and there is at most one such path, because two of them would
give a cycle of length at most $2s<D$. For the second statement, a cycle of length $\ell$ through $u$
can be traversed from $u$ in two directions, so there are at most $\half N_\ell(u,u)$ of them, and
$\ell>D>s$.
\end{proof}

Let $B$ be the non-backtracking matrix of $J$. Its rows and columns are indexed by the directed
edges $(i\to j)$ of $J$, two for each edge. The entry in row $(i\to j)$ and column $(j'\to k)$ is
$1$ if $j'=j$ and $k\ne i$, and it is $0$ otherwise.

\begin{lemma}\label{lem:spectrum}
For every even $\ell\ge2$,
\[
  \tr B^{\ell}\;\ge\;\bigl(2-d_J/v\bigr)^{\ell}-3v\,2^{\ell/2}.
\]
\end{lemma}

\begin{proof}
Let $A$ be the adjacency matrix of $J$ and $\Delta$ the diagonal matrix of its degrees, and put
$M(\lambda)=\lambda^{2}I-\lambda A+\Delta-I$.

(1) \emph{If $M(\lambda)g=0$ for a real $\lambda>1$ and a vector $g\ne0$, then $\lambda$ is an
eigenvalue of $B$.} A direct computation shows that $f(i\to j)=\lambda g(j)-g(i)$ satisfies
$Bf=\lambda f$; and $f=0$ would give $g(i)=\lambda g(j)=\lambda^{2}g(i)$ along every edge, hence
$g=0$ (at an isolated vertex, $M(\lambda)g=0$ already says $(\lambda^{2}-1)g=0$).

(2) \emph{If $Bf=\lambda f$ with $f\ne0$ and $\lambda^{2}\ne1$, then $M(\lambda)g=0$ for
$g(j)=\sum_{k\sim j}f(j\to k)$, and $g\ne0$.} The eigenvalue equation at $(i\to j)$ and at
$(j\to i)$ reads $\lambda f(i\to j)=g(j)-f(j\to i)$ and $\lambda f(j\to i)=g(i)-f(i\to j)$.
Eliminating $f(j\to i)$ gives $(\lambda^{2}-1)f(i\to j)=\lambda g(j)-g(i)$. This shows that
$g\ne0$. Summing it over the neighbours $j$ of $i$ gives $M(\lambda)g=0$.

(3) \emph{$B$ has a real eigenvalue $\lambda_1\ge2-d_J/v$.} The matrix
$M(2)=2(\Delta-A)+(3I-\Delta)$ is positive semidefinite, while
$\mathbf{1}^{\!\top}M(\lambda)\mathbf{1}=(\lambda-1)(v\lambda-2v+d_J)$ vanishes at
$\lambda^{*}=2-d_J/v>1$. The smallest eigenvalue of the symmetric matrix $M(\lambda)$ is continuous
in $\lambda$, at most $0$ at $\lambda^{*}$ and at least $0$ at $2$. So $M(\lambda_1)$ is singular for
some $\lambda_1\in[\lambda^{*},2]$, and (1) applies.

(4) \emph{Every non-real eigenvalue $\lambda$ of $B$ satisfies $|\lambda|\le\sqrt2$.} With $g$ as
in~(2), $g^{*}M(\lambda)g=0$ is a quadratic equation $a\lambda^{2}-b\lambda+c'=0$ with real
coefficients $a=g^{*}g>0$, $b=g^{*}Ag$ and $c'=g^{*}(\Delta-I)g\le2a$. Its roots are $\lambda$ and
$\bar\lambda$, so $|\lambda|^{2}=c'/a\le2$.

Now $\tr B^{\ell}$ is the sum of the $\ell$th powers of the $2|E(J)|\le3v$ eigenvalues of $B$. For
even $\ell$ the real eigenvalues contribute nonnegative terms, one of which is at least
$(2-d_J/v)^{\ell}$, and the non-real ones contribute at least $-3v\,2^{\ell/2}$ in total.
\end{proof}

The argument uses the same matrix-pencil proof as the existing manuscript.
It needs no expansion assumption.


\begin{lemma}[Repeated closed walks]\label{lem:polynomial-repeated-walks}
Let $\mathcal C$ contain the cycles of $J$ with length at most $L$ and put $\mu=\sum_{C\in\mathcal C}2^{-|C|}$.
Then
\[
 \mu\le T:=\sum_{\ell\le L}\frac{\operatorname{tr}B^\ell}{2\ell2^\ell}
 \le\left(1+\tfrac32L^2 2^{-s}\right)\mu.
\]
\end{lemma}
\begin{proof}
$\tr B^{\ell}$ is the number of closed walks $u_0u_1\cdots u_\ell=u_0$ in $J$ with
$u_{i+1}\ne u_{i-1}$ for all $i$ modulo $\ell$. A cycle of length $\ell$ gives $2\ell$ of them; let
$n_\ell$ be the number of the others. Then
\[
  T:=\sum_{\ell\le L}\frac{\tr B^{\ell}}{2\ell\,2^{\ell}}=\mu+\sum_{\ell\le L}\frac{n_\ell}{2\ell\,2^{\ell}} .
\]
Consider a walk counted by $n_\ell$, and let $j$ be the least index with
$u_j\in\{u_0,\dots,u_{j-1}\}$, say $u_j=u_i$. Then $j<\ell$, the walk $C=u_iu_{i+1}\cdots u_j$ is a
cycle, of length $q$ say, and $u_ju_{j+1}\cdots u_\ell u_1\cdots u_i$ is a non-backtracking closed
walk of length $\ell-q>0$ from $u_i$ to itself. It contains a cycle, so $\ell-q>D>s$. The original
walk is determined by $C$, by the vertex $u_i$ and a direction on $C$ ($2q$ choices), by the
complementary walk (at most $3\cdot2^{\ell-q-s-1}$ choices, by Lemma~\ref{lem:walks}), and by the
position of $u_0$ on the latter (at most $\ell$ choices). Hence
\[
  \sum_{\ell\le L}\frac{n_\ell}{2\ell\,2^{\ell}}
  \;\le\;\sum_{\ell\le L}\ \sum_{C\in\mathcal{C}}\tfrac32\,|C|\,2^{-|C|-s}
  \;\le\;\tfrac32\,L^{2}\,2^{-s}\,\mu,
\]
as required.


\end{proof}

In particular, suppose $v=O(r)$, $x=\log_2r$, $L=\lfloor Ax\rfloor$, and $A>6$ is fixed.
If $Ld_J=o(v)$ and $L^2 2^{-s}=o(1)$, then
\[
 \mu\ge(1/4-o(1))\ln(A/6).
\]
Indeed, sum the even trace bound over $6x\le\ell\le L$.
Its normalized leading term is $1-o(1)$, and its nonreal error is $O(r^{-2})$.
The even harmonic sum is $(1/4-o(1))\ln(A/6)$.
The repeated-walk error is $o(\mu)$.

\subsection{The uniform bound on links}
Suppose $\Gamma$ is connected and subcubic, $\Gamma[P]$ is connected, and every vertex outside $P$ has degree three.
Let $\mathcal Q$ be a maximal disjoint family of cycles of length at most $D$ in $\Gamma-P$.
Write $S$ for its vertex union and $J=\Gamma-(P\cup S)$.
Assume the probability of a forest outside $P$ exceeds $1/2+\varepsilon$.
For disjoint cycles $C,C'$ of $J$ with lengths at most $L$, and an integer $t$ with $2t+D/2\le D$,
\begin{equation}\label{eq:polynomial-foundation-links}
 \ell(C,C')\le1+\frac{2L}{t}+8t2^t\ln(2/\varepsilon)+(8L/D)^2.
\end{equation}

To prove this, let $U=V(C)\cup V(C')$.
At most one component of $\Gamma-U$ meets $P$.
At most $2L/t$ other components have more than $t$ boundary edges.
The exponential cut estimate bounds the number of cyclic components with at most $t$ boundary edges by $8t2^t\ln(2/\varepsilon)$.

A remaining tree component $Y$ has $|Y|=|\delta Y|-2\le t-2$.
It contains no chosen short cycle, since every chosen cycle avoids $U$ and is connected.
Thus $Y\subseteq J$.
If $Y$ is adjacent to both cycles, it supplies a connecting path of length at most $t$.
Direct edges between the cycles are also such paths.
Subcubic degree makes these paths vertex-disjoint, including their endpoints.
Divide each cycle into at most $8L/D$ arcs with at most $D/4$ consecutive vertices.
Two connecting paths with the same pair of endpoint arcs would produce a cycle of length at most $2t+D/2\le D$ in $J$.
There is therefore at most one path per pair of arcs.
This proves \eqref{eq:polynomial-foundation-links}.

\subsection{Component variables and exceptional pairs}
For a cycle $C\subseteq J$, put $a_C=2^{-|C|}$ and let $I_C$ mean that $C$ is a component of the uniform even subgraph $\xi$ of $\Gamma$.
Set $p_C=\Pr(I_C)$ and $X_C=(a_C/p_C)\mathbf1_{I_C}$.
The condition $E(C)\subseteq\xi$ is feasible and imposes at most $|C|$ binary linear equations.
Hence $p_C\ge a_C$, so $\mathbb EX_C=a_C$ and $\mathbb EX_C^2\le a_C$.
The link formula gives
\[
 \mathbb E(X_CX_{C'})=a_Ca_{C'}2^{\ell(C,C')-1}
\]
for disjoint cycles.
Distinct intersecting cycles cannot both be components.

Fix $C$ and choose a maximal disjoint family of partners with at least three links to $C$.
The exceptional-partner lemma bounds its size by $L$.
Its vertex union has size at most $L^2$ and meets every exceptional partner.
The local walk bound therefore bounds their total weight by $L^3 2^{-s}$.
If all link counts are at most $b$, their total normalized joint contribution is at most
\[
 a_C2^bL^3 2^{-s}.
\]
These estimates apply to each of the two graph deletions in the next section.
